\begin{pastexam}[label={px:gk3-104-ds-03}]{高考三級資訊處理 104 年 第三題}
（申論；20 分）

給定一個權重圖（weighted graph），$G=(V,E,w)$，其中每個邊（edge）$e$ 的權重 $w(e)$ 都是正整數，為了簡單，假設 $V=\{1,2,\ldots,n\}$。任意點 $v$ 與起始點 $s$ 的距離可以用一個矩陣 $d[1..n]$ 來表示。（每小題 10 分，共 20 分）
\begin{itemize}[leftmargin=2.2em]
\item[㈠] 設計一個只需 $O(n)$ 空間的方法來記錄從 $s$ 出發，到達每個點的最短路徑。
\item[㈡] 說明計算與印出從起始點 $s$ 到任意點 $t\in V$ 的最短路徑的演算法。（解此小題時可參考 Dijkstra 或其他演算法來設計，且不須將 Dijkstra 或別的演算法做詳細的描述。）
\end{itemize}
\end{pastexam}

\begin{pastsoln}
題目的「矩陣 $d[1..n]$」是一維陣列。 想法是資結書第~25.1~節的\textbf{前驅陣列}（最短路徑樹）。

\textbf{㈠ 前驅陣列}。 另開一個陣列 $\pi[1..n]$：$\pi[v]$ 存「從 $s$ 到 $v$ 的某一條最短路徑上，$v$ 的前一個
點」；$\pi[s] = 0$，到不了的點也存 $0$（$0$ 不是合法的點，當成 $\NIL$）。 只用 $n$ 個整數，是 $O(n)$ 空間。

為什麼一個前驅就夠：最短路徑的子路徑還是最短路徑（資結書引理~25.2）。 若 $s \leadsto u \to v$ 是到 $v$ 的
最短路徑，前半段 $s \leadsto u$ 就是一條到 $u$ 的最短路徑，所以「到 $v$ 的路徑」可以寫成「到 $\pi[v]$ 的
路徑再走一條邊」，一路往回接到 $s$。 所有邊 $(\pi[v], v)$ 合起來是一棵以 $s$ 為根的樹（最短路徑樹，資結書
定理~25.5），樹上 $s$ 到 $v$ 的路徑就是一條最短路徑。 雖然最短路徑可能有 $\Theta(n)$ 個點、共 $n$ 條，
全部存下來要 $\Theta(n^2)$ 空間，但它們共用前綴，存成樹只要 $O(n)$。

\textbf{㈡ 計算與印出}。 權重都是正數，用 Dijkstra 演算法（資結書第~25.2~節）。 計算時在每次鬆弛成功的地方
多記一筆前驅：
\begin{pseudo}[unbreakable]{鬆弛時順便記錄前驅（其餘照 Dijkstra）}
\Procedure{Relax}{$u, v$}
    \If{$d[u] + w(u, v) < d[v]$}
        \State $d[v] \gets d[u] + w(u, v)$
        \State $\pi[v] \gets u$ \Comment{目前到 $v$ 的最好路徑是經過 $u$}
    \EndIf
\EndProcedure
\end{pseudo}
初始時 $d[s] = 0$、其餘 $d[v] = \infty$、所有 $\pi[v] = 0$。 Dijkstra 結束時 $d[v]$ 是最短距離；每個點的
$\pi$ 最後一次被改，是在 $d[v]$ 被改成最終值的那一次，那時的 $u$ 已經有最終距離且 $d[u] + w(u, v) = d[v]$，
所以 $(\pi[v], v)$ 確實是某條最短路徑的最後一條邊。 多記 $\pi$ 不改變 Dijkstra 的時間（例如用二元堆積是
$O((n + m) \log n)$，$m = |E|$）。

印出 $s$ 到 $t$ 的路徑：從 $t$ 沿著 $\pi$ 往回走到 $s$，得到的是倒過來的順序，所以用遞迴（或先推進堆疊再
依序彈出）把它轉正：
\begin{pseudo}[unbreakable]{印出 $s$ 到 $t$ 的最短路徑}
\Procedure{Print-Path}{$s, t$}
    \If{$t = s$}
        \State 印出 $s$
    \ElsIf{$\pi[t] = 0$}
        \State 印出「$s$ 到不了 $t$」
    \Else
        \State \Call{Print-Path}{$s, \pi[t]$}
        \State 印出 $t$
    \EndIf
\EndProcedure
\end{pseudo}
每一層遞迴往回走一條邊，路徑最多 $n - 1$ 條邊，所以印出的時間是 $O(n)$（正比於路徑長度），遞迴深度也最多
$n$。 路徑長度就是 $d[t]$。

例：$n = 4$、$s = 1$，邊 $1 \to 2$（$4$）、$1 \to 3$（$1$）、$3 \to 2$（$2$）、$2 \to 4$（$5$）。 Dijkstra 依序
確定 $1$（$0$）、$3$（$1$）、$2$（$3$，由 $3$ 鬆弛，$\pi[2]$ 從 $1$ 改成 $3$）、$4$（$8$）。 結果
$\pi = [0, 3, 1, 2]$、$d = [0, 3, 1, 8]$。 印 $t = 4$：$\pi[4] = 2$、$\pi[2] = 3$、$\pi[3] = 1 = s$，
遞迴展開後印出 $1, 3, 2, 4$，長度 $d[4] = 8$。
\end{pastsoln}
